% !TeX root = ../main.tex
% Fonte: Lezione 10, p. 2.

\section{Prodotto di spazi vettoriali}

\subsection{Operazioni nel prodotto}

Siano $V,W$ spazi vettoriali sullo stesso campo $K$. Sul prodotto cartesiano
\[
  V\times W=\{(v,w):v\in V,\ w\in W\}
\]
si definiscono le operazioni componente per componente:
\begin{align*}
  (v_1,w_1)+(v_2,w_2)&=(v_1+v_2,w_1+w_2),\\
  \lambda(v,w)&=(\lambda v,\lambda w),\qquad\lambda\in K.
\end{align*}
Entrambe hanno risultato in $V\times W$, per la chiusura delle operazioni nei fattori.

\begin{proposizione}
Con queste operazioni $V\times W$ è uno spazio vettoriale su $K$.
\end{proposizione}
\begin{proof}
Poniamo $p=(v,w)$, $q=(v',w')$ e $t=(v'',w'')$.

\textbf{Struttura additiva.} Commutatività e associatività si verificano nelle componenti:
\begin{align*}
  p+q&=(v+v',w+w')=(v'+v,w'+w)=q+p,\\
  (p+q)+t&=((v+v')+v'',(w+w')+w'')\\
    &=(v+(v'+v''),w+(w'+w''))=p+(q+t).
\end{align*}
Il neutro è $(0_V,0_W)$ e l'opposto di $(v,w)$ è $(-v,-w)$, come si verifica sommando per componenti.

\textbf{Prodotto per scalari.} Per $\lambda,\mu\in K$, gli assiomi nei due fattori danno
\begin{align*}
  \lambda(p+q)&=(\lambda(v+v'),\lambda(w+w'))=\lambda p+\lambda q,\\
  (\lambda+\mu)p&=((\lambda+\mu)v,(\lambda+\mu)w)=\lambda p+\mu p,\\
  \lambda(\mu p)&=(\lambda(\mu v),\lambda(\mu w))\\
    &=((\lambda\mu)v,(\lambda\mu)w)=(\lambda\mu)p,\\
  1_Kp&=(1_Kv,1_Kw)=(v,w)=p.
\end{align*}
Sono soddisfatti tutti gli assiomi.
\end{proof}

\subsection{Base e dimensione del prodotto}

\begin{proposizione}
Se $(v_1,\ldots,v_r)$ è una base di $V$ e $(w_1,\ldots,w_s)$ è una base di $W$, allora
\[
  (v_1,0_W),\ldots,(v_r,0_W),\quad(0_V,w_1),\ldots,(0_V,w_s)
\]
è una base di $V\times W$. In particolare,
\[
  \dim(V\times W)=\dim V+\dim W.
\]
\end{proposizione}
\begin{proof}
\textbf{Generazione.} Per $(v,w)\in V\times W$, scriviamo $v=\sum_i a_iv_i$ e $w=\sum_j b_jw_j$. Allora
\[
  (v,w)=\sum_{i=1}^r a_i(v_i,0_W)+\sum_{j=1}^s b_j(0_V,w_j).
\]

\textbf{Indipendenza.} Una combinazione nulla della famiglia proposta dà
\[
  \sum_{i=1}^r a_i(v_i,0_W)+\sum_{j=1}^s b_j(0_V,w_j)=(0_V,0_W)
  \quad\Longrightarrow\quad
  \begin{cases}
    \sum_{i=1}^r a_iv_i=0_V,\\
    \sum_{j=1}^s b_jw_j=0_W.
  \end{cases}
\]
L'indipendenza delle basi dei fattori annulla tutti gli $a_i$ e i $b_j$. La famiglia è quindi una base, con $r+s$ vettori: segue la formula della dimensione.
\end{proof}
